% Faithful official CCC2022 article16 Definition39, Theorem40 and Theorem43 with complete author proofs.
\section*{4.4 Case of PH (main promise-PH simulation)}
It is observed in [7, Corollary 17] that, for every $\ell\geq1$, $\Sigma_\ell^{\mathrm P}=\operatorname{lin}\Sigma_\ell^{\mathrm P}$, under the assumption that $\mathrm{Dist}\Sigma_{\ell+1}^{\mathrm P}\subseteq\mathrm{AvgP}$. We will show a similar result for the class $\operatorname{promise}\text{-}\Sigma_\ell^{\mathrm{BP}}$ defined below.
\paragraph{Definition 39 ($\operatorname{promise}\text{-}\Sigma_\ell^{\mathrm{BP}}$ and $\operatorname{promise}\text{-}\operatorname{lin}\Sigma_\ell^{\mathrm{BP}}$).} A promise problem $\Pi=(\Pi_Y,\Pi_N)\in\operatorname{promise}\text{-}\Sigma_\ell^{\mathrm{BP}}$ for $\ell\geq0$, if there is a $\Sigma_\ell\circ\mathrm{BP}\circ\mathrm P$ formula
\[\varphi(x)=\exists y_1\forall y_2\ldots Qy_\ell\ \mathrm{BP}_r\ R(x,y_1,y_2,\ldots,y_\ell,r),\]
for a polytime predicate $R$ and all string variables $y_1,\ldots,y_\ell,r$ of length polynomial in $|x|$, such that, for all $x\in\{0,1\}^n$,
\begin{align*}
x\in\Pi_Y&\Longrightarrow\exists y_1\forall y_2\ldots Qy_\ell\ \mathrm{BP}_{(2/3)}r\ R(x,y_1,y_2,\ldots,y_\ell,r),\text{ and}\\
x\in\Pi_N&\Longrightarrow\forall y_1\exists y_2\ldots\overline Qy_\ell\ \mathrm{BP}_{(2/3)}r\ \neg R(x,y_1,y_2,\ldots,y_\ell,r),
\end{align*}
where $\mathrm{BP}_{(\delta)}r\ R(-,r)$ means that $\Pr_r[R(-,r)]\geq\delta$. The constant $2/3$ in the definition above can be changed to $1-2^{-\operatorname{poly}(n)}$ by standard success amplification techniques (even if we start with any pair $\delta_Y>1-\delta_N+\varepsilon$ of thresholds, for some $\varepsilon\geq1/\operatorname{poly}(n)$, where $\delta_Y$ is the probability $R(-,r)$ accepts, and $\delta_N$ the probability $\neg R(-,r)$ accepts).
The promise class $\operatorname{promise}\text{-}\operatorname{lin}\Sigma_\ell^{\mathrm{BP}}$ is defined the same way as above, with an extra condition that, for some constant $a>0$, $|y_i|\leq an$ for all $1\leq i\leq\ell$ (but the randomness $r$ may still be of any $\operatorname{poly}(n)$ length).
\paragraph{Theorem 40 (Witness Compression for PH).} For any $\ell\geq0$, suppose $\mathrm{Dist}\Sigma_{\ell+2}^{\mathrm P}\subseteq\mathrm{AvgBPP}$. Then
\[\operatorname{promise}\text{-}\Sigma_\ell^{\mathrm{BP}}=\operatorname{promise}\text{-}\operatorname{lin}\Sigma_\ell^{\mathrm{BP}}.\]
\paragraph{Proof.} The proof is by induction on $\ell$. The base case of $\ell=0$ is trivially true.
For any $\ell>0$, suppose the claim is true for $\ell-1$. Consider an arbitrary promise problem $\Pi=(\Pi_Y,\Pi_N)\in\operatorname{promise}\text{-}\Sigma_\ell^{\mathrm{BP}}$. That is, for every $x\in\{0,1\}^n$,
\begin{align*}
x\in\Pi_Y&\Longrightarrow\exists y\ \Pi_{\ell-1}z\ \mathrm{BP}_{(1-2^{-n})}r\ R(x,y,z,r),\\
x\in\Pi_N&\Longrightarrow\forall y\ \Sigma_{\ell-1}z\ \mathrm{BP}_{(1-2^{-n})}r\ \neg R(x,y,z,r),
\end{align*}
for some polytime predicate $R$, where $z$ is a sequence of $\ell-1$ strings $z_1,\ldots,z_{\ell-1}$ under the corresponding $\ell-1$ quantifiers of the $\Pi_{\ell-1}$ pattern, and $|y|,|r|,|z_i|\leq q(n)$, for some polynomial $q$, for all $1\leq i\leq\ell-1$.
As in the proof of the inclusion $\mathrm{BPP}\subseteq\Sigma_2^{\mathrm P}\cap\Pi_2^{\mathrm P}$ [22], we replace the BP quantifier by the $\exists\forall$ pattern if $\ell$ is odd, or by the $\forall\exists$ pattern if $\ell$ is even. We get that there exists a polytime predicate $R'$ such that, for every $x\in\{0,1\}^n$,
\begin{align*}
x\in\Pi_Y&\Longrightarrow\exists y\ \Pi_{\ell-1}z\ Q_0z'\ Q_1z''\ R'(x,y,z,z',z''),\\
x\in\Pi_N&\Longrightarrow\forall y\ \Sigma_{\ell-1}z\ Q_1z'\ Q_0z''\ \neg R'(x,y,z,z',z''),
\end{align*}
where $z',z''$ have the lengths polynomial in $n$ and $Q_0Q_1=\exists\forall$ if $\ell$ is odd, and $Q_0Q_1=\forall\exists$ if $\ell$ is even.
Note that the above transformation shows that our promise problem $\Pi\in\operatorname{promise}\text{-}\Sigma_{\ell+1}^{\mathrm P}$. For any given $x\in\Pi_Y$, let $y_x$ be the lexicographically first string $y$ such that
\[\Pi_{\ell-1}z\ Q_0z'\ Q_1z''\ R'(x,y,z,z',z'').\tag{23}\]
\paragraph{Claim 41.} For $x$ and $y_x$ as above,
\[\Pi_{\ell-1}z\ \mathrm{BP}_{(1/\operatorname{poly}(n))}r\ R(x,y_x,z,r).\]
\paragraph{Proof.} Immediate from the properties of Lautemann's proof of $\mathrm{BPP}\subseteq\Sigma_2^{\mathrm P}$ [22]. Recall that in that proof, the subformula $\mathrm{BP}_r\ R(-,r)$ is replaced by
\[\exists u_1\ldots u_k\ \forall r\ \bigvee_{i=1}^kR(-,r\oplus u_i),\]
for some $k\leq\operatorname{poly}(|r|)$, where $|u_i|=|r|$, for all $1\leq i\leq k$, and $\oplus$ is the bit-wise XOR. If $\Pr_r[R(-,r)]<1/k$, then such a collection of $u_1,\ldots,u_k$ cannot exist. Indeed, each ``shifted'' predicate $R(-,r\oplus u_i)$, for $1\leq i\leq k$, accepts less than $1/k$ fraction of $r\in\{0,1\}^{|r|}$, and hence, by the union bound, the OR of any $k$ such shifts accepts less than $k\cdot(1/k)=1$ fraction of $r\in\{0,1\}^{|r|}$.\hfill$\triangleleft$
On the other hand, for every $x\in\Pi_N$ and every $y$, we have
\[\Sigma_{\ell-1}z\ \mathrm{BP}_{(1-2^{-n})}r\ \neg R(x,y,z,r).\]
Using Claim 41, by standard success amplification, for a random string $r'$ of length $q'(n)$, for some polynomial $q'$, and a polytime predicate $R^{\mathrm{boost}}$, we get for every $x\in\{0,1\}^n$ and $y_x$ as defined above,
\begin{align*}
x\in\Pi_Y&\Longrightarrow\Pi_{\ell-1}z\ \mathrm{BP}_{(1-2^{-n})}r'\ R^{\mathrm{boost}}(x,y_x,z,r'),\tag{24}\\
x\in\Pi_N&\Longrightarrow\forall y\ \Sigma_{\ell-1}z\ \mathrm{BP}_{(1-2^{-n})}r'\ \neg R^{\mathrm{boost}}(x,y,z,r').\tag{25}
\end{align*}
\paragraph{Claim 42.} For some universal constant $a>0$ and some polynomial $p'$, for every $n$-bit $x\in\Pi_Y$ with a witness $y_x$ as defined above, we have
\[\mathrm{rK}^{p'(n)}(y_x\mid x)\leq an.\]
\paragraph{Proof.} Using Equation (23), we get by Item 2 of Lemma 34 (since $\mathrm{Dist}\Sigma_3^{\mathrm P}\subseteq\mathrm{AvgBPP}$), that, for some polynomial $p$ and all large enough $t$,
\[\mathrm{rK}^{p(t)}(y_x\mid x)\leq\mathrm{rK}^t(x)-\mathrm{rK}^{p(t)}(x)+2\log p(t).\]
Since, for any $t$, $\mathrm{rK}^t(x)\leq O(n)$, the claim follows.\hfill$\triangleleft$
Consider the formula
\[\varphi(x)=\exists y'\ (|y'|\leq an)\ \Pi_{\ell-1}z\ \mathrm{BP}_w\ R''(x,y',z,w),\tag{26}\]
where the predicate $R''$ is computed by the following polytime algorithm $A$:
\begin{itemize}
\item On input $x,y',z$, and a string $w=w'r'$, where $w'$ is of length $p'(n)$, and $r'$ is of length $q'(n)$, first run the universal machine $U(y',x,w')$ for $p'(n)$ steps, getting an output $y$. Then run $R^{\mathrm{boost}}(x,y,z,r')$, accepting iff $R^{\mathrm{boost}}$ accepts.
\end{itemize}
We claim that Equation (26) is a correct $\Sigma_\ell^{\mathrm{BP}}$ formula for the promise problem $\Pi$. Indeed, for any $x\in\Pi_Y$, by Claim 42, there is a string $y'$ that is decompressed to $y_x$ with high probability over random strings $w'$ (say, with probability at least $3/4$). By Equation (24), $R^{\mathrm{boost}}(x,U(y',x,w'),z,r')$ accepts with probability at least $(3/4)(1-2^{-n})>2/3$ over strings $w=w'r'$. On the other hand, for any $x\in\Pi_N$, Equation (25) implies that
\[\forall y'\ \Sigma_{\ell-1}z\ \mathrm{BP}_{(1-2^{-n})}w\ \neg R''(x,y',z,w).\]
Next, using Equation (26), define a new promise problem $\Pi'=(\Pi'_Y,\Pi'_N)\in\operatorname{promise}\text{-}\Pi_{\ell-1}^{\mathrm{BP}}$:
\begin{align*}
\Pi'_Y&=\{(x,y')\mid|y'|\leq a|x|,\ \Pi_{\ell-1}z\ \mathrm{BP}_w\ R''(x,y',z,w)\},\\
\Pi'_N&=\{(x,y')\mid|y'|\leq a|x|,\ \Sigma_{\ell-1}z\ \mathrm{BP}_w\ \neg R''(x,y',z,w)\}.
\end{align*}
By the Inductive Hypothesis, $\Pi'\in\operatorname{promise}\text{-}\operatorname{lin}\Pi_{\ell-1}^{\mathrm{BP}}$, with all quantified binary strings of length $O(|x|+|y'|)\leq O(n)$. It follows that $\Pi\in\operatorname{promise}\text{-}\operatorname{lin}\Sigma_\ell^{\mathrm{BP}}$.\hfill$\triangleleft$
\paragraph{Theorem 43.} For any $\ell\geq0$,
\[\mathrm{Dist}\Sigma_{\ell+2}^{\mathrm P}\subseteq\mathrm{AvgBPP}\Longrightarrow\operatorname{promise}\text{-}\Sigma_\ell^{\mathrm P}\subseteq\operatorname{promise}\text{-}\mathrm{BPTIME}[2^{O(n/\log n)}].\]
\paragraph{Proof.} The proof is by induction on $\ell$. The base case of $\ell=0$ is trivially true. Consider any $\ell\geq1$, and assume the claim for $\ell-1$. Let $\Pi=(\Pi_Y,\Pi_N)\in\operatorname{promise}\text{-}\Sigma_\ell^{\mathrm P}$ be arbitrary. By Theorem 40, $\Pi\in\operatorname{promise}\text{-}\operatorname{lin}\Sigma_\ell^{\mathrm{BP}}$ via some formula
\[\varphi(x)=\exists y\ \Pi_{\ell-1}z\ \mathrm{BP}_r\ R(x,y,z,r),\]
with a polytime $R$, linearly-bounded $y$ and $z=(z_1,\ldots,z_{\ell-1})$ (under the $\Pi_{\ell-1}$ quantifiers), and a poly-bounded $r$.
Arguing as in the proof of Theorem 40 (see the proof of Claim 42), but using Lemma 31 (for the case of pK) to bound the computational depth of an input $x$ by $O(n/\log n)$ with a time bound $t\leq2^{O(n/\log n)}$, we get that, for some universal constant $a>0$, every $x\in\Pi_Y$ has a witness $y_x$ with
\[\mathrm{pK}^{2^{an/\log n}}(y_x\mid x)\leq an/\log n.\tag{27}\]
Applying success amplification to $R$, we get a new polytime predicate $R^{\mathrm{boost}}$ with poly-bounded randomness $r'$ such that we get for every $x\in\{0,1\}^n$ and $y_x$ as defined above,
\begin{align*}
x\in\Pi_Y&\Longrightarrow\Pi_{\ell-1}z\ \mathrm{BP}_{(1-2^{-n})}r'\ R^{\mathrm{boost}}(x,y_x,z,r'),\tag{28}\\
x\in\Pi_N&\Longrightarrow\forall y\ \Sigma_{\ell-1}z\ \mathrm{BP}_{(1-2^{-n})}r'\ \neg R^{\mathrm{boost}}(x,y,z,r').\tag{29}
\end{align*}
Define a promise problem $\Pi'=(\Pi'_Y,\Pi'_N)\in\operatorname{promise}\text{-}\Pi_{\ell-1}^{\mathrm{BP}}$:
\begin{align*}
\Pi'_Y&=\{(x,y)\mid\Pi_{\ell-1}z\ \mathrm{BP}_{r'}\ R^{\mathrm{boost}}(x,y,z,r')\},\\
\Pi'_N&=\{(x,y)\mid\Sigma_{\ell-1}z\ \mathrm{BP}_{r'}\ \neg R^{\mathrm{boost}}(x,y,z,r')\}.
\end{align*}
By the Inductive Hypothesis, $\Pi'\in\operatorname{promise}\text{-}\mathrm{BPTIME}[2^{O(n/\log n)}]$, where $|x|=n$, via some randomized algorithm $A$, with error probability at most $2^{-n}$.
We now solve $\Pi$ with the following randomized algorithm $B$:
\begin{itemize}
\item Given $x\in\{0,1\}^n$, choose a random string $w$ of length $2^{an/\log n}$, and then enumerate over all machines $M\in\{0,1\}^{an/\log n}$, and run $M(w,x)$ for $2^{an/\log n}$ steps, getting a candidate witness $y$ of length $O(n)$. Run $A(x,y)$, accepting iff $A$ accepts.
\end{itemize}
The time complexity of the described algorithm $B$ is clearly $2^{O(n/\log n)}$. For correctness, for every $x\in\Pi_Y$, with probability at least $3/4$ over $w$, our algorithm will consider a short description of a TM $M$ such that $M(w,x)=y_x$. Since $(x,y_x)\in\Pi'_Y$, we get that $A(x,y_x)$ will accept with probability at least $1-2^{-n}$ over its internal randomness. Hence, the algorithm $B$ accepts $x$ with probability at least $(3/4)(1-2^{-n})>2/3$. On the other hand, for every $x\in\Pi_N$, for every $w$ and all $M$, the string $y=M(w,x)$ is such that $(x,y)\in\Pi'_N$, and so $A(x,y)$ accepts with probability at most $2^{-n}$. It follows that the probability that $B(x)$ accepts in this case is
\[\mathbb E_{w,A}[\exists v\in\{0,1\}^{an/\log n}:A(x,U(v,x,w))]\leq2^{an/\log n}\cdot\mathbb E_w(2^{-n})\leq2^{-n/2},\]
for all sufficiently large $n$.\hfill$\triangleleft$
